% ======================================================================
% Route: the lifted bilinear system and its RLT hierarchy over the choice
% variables.  Labels prefixed cl:.  Paste-ready, amsthm style of frontier.tex.
% ======================================================================

\subsection{Lifting the choice: a hierarchy over the transport polytope}
\label{sec:choicelift}

\Cref{def:transport} keeps, of each controlled row of $Tx=x$, only its convex
half. The missing halves are not linear, but they become linear once the choice
itself is made a variable. This subsection carries that lift out, identifies
the resulting hierarchy's zeroth member with $Q(G)$, computes the level at
which it collapses, and exhibits the instance on which its first level is not
exact.

\begin{definition}[The choice lift]\label{cl:lift}
Let $G$ be an SSG with $C=\Vmax\cup\Vmin$. The \emph{lifted system} $L(G)$ is
the set of pairs $(x,y)\in\mathbb R^{V}\times[0,1]^{C}$ with
\[
 \begin{aligned}
  &x(t_0)=0,\quad x(t_1)=1,\quad 0\le x\le 1,\\
  &x(v)=\tfrac12\bigl(x(v^{(0)})+x(v^{(1)})\bigr) &&(v\in\Vavg),\\
  &x(v)\ge x(v^{(i)})\ \ (i=0,1) &&(v\in\Vmax),\\
  &x(v)\le x(v^{(i)})\ \ (i=0,1) &&(v\in\Vmin),\\
  &x(v)=y_v\,x(v^{(0)})+(1-y_v)\,x(v^{(1)}) &&(v\in C).
 \end{aligned}
\]
The last family is bilinear; all the others are the rows of $Q(G)$
(\Cref{def:transport}). Write $L_{\mathrm{int}}(G):=L(G)\cap\{y\in\{0,1\}^{C}\}$.
\end{definition}

\begin{proposition}[The lift is exact, and already convex]\label{cl:exact-lift}
Let $G$ be stopping. Then
\begin{enumerate}[label=(\alph*),leftmargin=2.2em]
\item $(x,y)\in L(G)$ implies $x=w^{\ast}$; no integrality is needed;
\item $L(G)=\{w^{\ast}\}\times\operatorname{conv}(\Opt)$, where
      $\Opt\subseteq\{0,1\}^{C}$ is the optimal subcube of
      \Cref{thm:opt-subcube} and $\operatorname{conv}(\Opt)$ is the face of
      $[0,1]^{C}$ it spans;
\item hence $L(G)=\operatorname{conv}\bigl(L_{\mathrm{int}}(G)\bigr)$: the
      lifted feasible set is its own convex hull, that hull is cut out by the
      $N$ equalities $x=w^{\ast}$ together with at most $|C|$ facets
      $y_v\ge0$, $y_v\le1$ of the cube, and it has no disjunctive structure of
      $2^{|\Vmin|}$ pieces at all.
\end{enumerate}
\end{proposition}

\begin{proof}
(a) At $v\in\Vmax$ the two order rows give $x(v)\ge\max_ix(v^{(i)})$ while the
bilinear row writes $x(v)$ as a convex combination of the two, so
$x(v)\le\max_ix(v^{(i)})$; hence $x(v)=\max_ix(v^{(i)})$. Dually
$x(v)=\min_ix(v^{(i)})$ at $v\in\Vmin$, and the average rows and the sink
values are those of $T$. So $Tx=x$ and $x=w^{\ast}$ by
\Cref{thm:contraction}.

(b) By (a) the bilinear row at $v\in\Vmax$ reads
$w^{\ast}(v)=y_vw^{\ast}(v^{(0)})+(1-y_v)w^{\ast}(v^{(1)})$ with
$w^{\ast}(v)=\max_iw^{\ast}(v^{(i)})$, so $y_v$ is supported on the argmax:
$y_v=1$ if $w^{\ast}(v^{(0)})>w^{\ast}(v^{(1)})$, $y_v=0$ in the opposite case,
and $y_v$ free in $[0,1]$ when the two are equal; dually at $\Vmin$. The set of
integral selections of these argmax/argmin sets is exactly $\Opt$, a subcube
whose free coordinates are the tied vertices (\Cref{thm:opt-subcube}), and the
product of $\{0\},\{1\},[0,1]$ over the coordinates is precisely the face it
spans. Conversely every such pair satisfies all the rows, $w^{\ast}$ lying in
$Q(G)$ by \Cref{thm:transport-sound}. (c) is immediate from (b).
\end{proof}

\begin{remark}\label{cl:hull-reading}
\Cref{cl:exact-lift}(c) closes the disjunctive-hull question in the vacuous
direction, and it is worth being exact about what it does and does not say.
Balas' description of the hull of a union of polyhedra would give, for the
system \emph{without} the order rows of $Q(G)$, a lift with one block of
variables per branch; with the order rows present there is nothing to
describe, because the order rows already annihilate every non-optimal branch:
$x_{\sigma\tau}\in Q(G)$ forces $Tx_{\sigma\tau}=x_{\sigma\tau}$. So the hull
is small and useless, and the difficulty is not that $\operatorname{conv}
(L_{\mathrm{int}})$ has many facets --- it has at most $|C|$ --- but that the
\emph{linearisations} below approach it slowly. Dropping the order rows
instead gives the profile hull $\operatorname{conv}\{x_{\sigma\tau}\}$, over
which no linear functional is $w^{\ast}$: optimising there computes
$\max_{\sigma\tau}$, not $\max_\sigma\min_\tau$.
\end{remark}

\begin{definition}[The level-$j$ relaxation]\label{cl:rlt}
Fix $j\ge0$. A \emph{level-$j$ point} of $G$ is a family
$\bigl(z^{S,T}\in\mathbb R,\;X^{S,T}\in\mathbb R^{V}\bigr)$ indexed by the
ordered pairs of disjoint sets $S,T\subseteq C$ with $|S|+|T|\le j$, such that
\begin{enumerate}[label=(\roman*),leftmargin=2.4em]
\item $z^{\emptyset,\emptyset}=1$;
\item (splitting) for $|S|+|T|<j$ and $v\notin S\cup T$,
      $z^{S,T}=z^{S\cup v,T}+z^{S,T\cup v}$ and
      $X^{S,T}=X^{S\cup v,T}+X^{S,T\cup v}$;
\item (homogenised transport rows) for every $(S,T)$:
      $X^{S,T}(t_0)=0$, $X^{S,T}(t_1)=z^{S,T}$,
      $0\le X^{S,T}\le z^{S,T}\mathbf 1$,
      $X^{S,T}(v)=\frac12(X^{S,T}(v^{(0)})+X^{S,T}(v^{(1)}))$ on $\Vavg$,
      $X^{S,T}(v)\ge X^{S,T}(v^{(i)})$ on $\Vmax$ and
      $X^{S,T}(v)\le X^{S,T}(v^{(i)})$ on $\Vmin$, for $i=0,1$;
\item (tightness) $X^{S,T}(v)=X^{S,T}(v^{(0)})$ for $v\in S$ and
      $X^{S,T}(v)=X^{S,T}(v^{(1)})$ for $v\in T$.
\end{enumerate}
Put $R_j(G):=\{X^{\emptyset,\emptyset}:\text{a level-}j\text{ point exists}\}$.
This is the projection of a linear programme with
$\bigl(\sum_{i\le j}\binom{|C|}{i}2^{i}\bigr)(N+1)$ variables and $O(N)$ rows
per index, hence of size $N^{O(j)}$ for fixed $j$, with every coefficient in
$\{0,\pm\frac12,\pm1\}$. It is the level-$j$ Reformulation--Linearisation lift
of \Cref{cl:lift} in the choice variables --- $z^{S,T}$ and $X^{S,T}$ stand for
$\mathbb E\prod_{v\in S}y_v\prod_{v\in T}(1-y_v)$ and for the same product
times $x$ --- but we take the displayed system as the definition, so that
nothing is imported.
\end{definition}

\begin{proposition}[Soundness, base, monotonicity]\label{cl:sound}
For every stopping $G$: $w^{\ast}\in R_j(G)$ for all $j$;
$R_0(G)=Q(G)$; $R_{j+1}(G)\subseteq R_j(G)$. Consequently
\[
  \min_{x\in R_j}x(v)\ \le\ w^{\ast}(v)\ \le\ \max_{x\in R_j}x(v)
  \qquad(v\in V),
\]
and both bounds are computable exactly in time $N^{O(j)}$.
\end{proposition}

\begin{proof}
Let $\sigma^{\ast}$ be an optimal profile and put $z^{S,T}:=1$ if
$\sigma^{\ast}(v)=0$ for all $v\in S$ and $\sigma^{\ast}(v)=1$ for all $v\in T$,
and $z^{S,T}:=0$ otherwise; $X^{S,T}:=z^{S,T}w^{\ast}$. Splitting holds because
exactly one of the two children of a pattern agrees with $\sigma^{\ast}$ at
$v$; (iii) holds because $w^{\ast}\in Q(G)$ (\Cref{thm:transport-sound}) and
the rows are homogeneous; (iv) holds because $\sigma^{\ast}$ is optimal, so
$w^{\ast}(v)=w^{\ast}(v^{(\sigma^{\ast}(v))})$, and is vacuous where $z^{S,T}=0$.
At $j=0$ the only index is $(\emptyset,\emptyset)$, (ii) and (iv) are empty and
(iii) is the definition of $Q(G)$. Monotonicity: restrict a level-$(j+1)$ point
to the indices of size at most $j$.
\end{proof}

\begin{lemma}[Level one is an intersection of hulls of faces]\label{cl:faces}
For $v\in C$ and $i\in\{0,1\}$ put
$F_v^{i}:=\{x\in Q(G):x(v)=x(v^{(i)})\}$. Each nonempty $F_v^{i}$ is a face of $Q(G)$
(the one on which the affine functional $x\mapsto x(v^{(i)})-x(v)$, of constant
sign on $Q(G)$, attains its extremum), and
\[
  R_1(G)\;=\;\bigcap_{v\in C}\operatorname{conv}\bigl(F_v^{0}\cup F_v^{1}\bigr).
\]
\end{lemma}

\begin{proof}
Fix $v$ and write $\lambda:=z^{\{v\},\emptyset}$, so $z^{\emptyset,\{v\}}
=1-\lambda$ by splitting, and both are $\ge0$ by (iii) applied to $X\ge0$ and
$X\le z\mathbf 1$ at $t_1$. If $\lambda>0$ then $X^{\{v\},\emptyset}/\lambda$
satisfies the rows of $Q(G)$ and, by (iv), lies in $F_v^{0}$; if $\lambda=0$
then (iii) forces $0\le X^{\{v\},\emptyset}\le0$, so
$X^{\{v\},\emptyset}=0$ and $x=X^{\emptyset,\{v\}}\in F_v^{1}$. In every case
$x\in\operatorname{conv}(F_v^{0}\cup F_v^{1})$. Conversely, from
$x=\lambda_v\pi^{v}+(1-\lambda_v)\kappa^{v}$ with $\pi^{v}\in F_v^{0}$,
$\kappa^{v}\in F_v^{1}$, set $z^{\{v\},\emptyset}:=\lambda_v$,
$X^{\{v\},\emptyset}:=\lambda_v\pi^{v}$, and dually; all of (i)--(iv) hold.
\end{proof}

\begin{definition}[Rigid and reversing vertices]\label{cl:rigid-def}
Call $v\in C$ \emph{rigid} (in $G$) if the affine functional
$A_v(x):=x(v^{(0)})-x(v^{(1)})$ has constant weak sign on $Q(G)$, i.e.\ if
$A_v\ge0$ throughout $Q(G)$ or $A_v\le0$ throughout $Q(G)$; otherwise $v$ is
\emph{reversing}. Write $\mathrm{Rev}(G)\subseteq C$ for the set of reversing vertices.
Rigidity is decided by $2|C|$ linear programmes over $Q(G)$.
\end{definition}

\begin{lemma}[A rigid vertex is pinned at level one]\label{cl:rigid-pin}
Let $v\in C$ be rigid and let $x\in\operatorname{conv}(F_v^{0}\cup F_v^{1})$.
Then $x(v)=\max_ix(v^{(i)})$ if $v\in\Vmax$, and $x(v)=\min_ix(v^{(i)})$ if
$v\in\Vmin$.
\end{lemma}

\begin{proof}
Say $v\in\Vmin$ and $A_v\ge0$ on $Q(G)$, so $x(v^{(1)})=\min_ix(v^{(i)})$
throughout $Q(G)$. If $\pi\in F_v^{0}$ then
$\pi(v)=\pi(v^{(0)})\ge\pi(v^{(1)})\ge\pi(v)$, the last by the order row, so
$\pi(v)=\pi(v^{(1)})$ and $\pi\in F_v^{1}$. Hence $F_v^{0}\subseteq F_v^{1}$
and $\operatorname{conv}(F_v^{0}\cup F_v^{1})=F_v^{1}$, which is convex; so
$x(v)=x(v^{(1)})=\min_ix(v^{(i)})$. The three remaining sign/side cases are
identical.
\end{proof}

\begin{theorem}[Rigid games are one linear programme]\label{cl:rigid}
If every controlled vertex of a stopping $G$ is rigid, then $R_1(G)=\{w^{\ast}\}$.
The value vector is therefore the unique point of an explicit linear
programme with $O(N|C|)$ variables and rows and coefficients in
$\{0,\pm\frac12,\pm1\}$, and the hypothesis is checked by $2|C|$ linear
programmes.
\end{theorem}

\begin{proof}
Let $x\in R_1(G)$. By \Cref{cl:faces} and \Cref{cl:rigid-pin}, $x$ satisfies
$x(v)=\max_i x(v^{(i)})$ at every Max vertex and $x(v)=\min_ix(v^{(i)})$ at
every Min vertex; the average rows and sinks are part of $Q(G)$. So $Tx=x$ and
$x=w^{\ast}$ by \Cref{thm:contraction}; and $w^{\ast}\in R_1$ by
\Cref{cl:sound}.
\end{proof}

\begin{theorem}[The collapse level]\label{cl:collapse}
Let $G$ be stopping, $k:=|\Vmin|$, $m:=|\Vmax|$,
$\rho_{\min}:=|\mathrm{Rev}(G)\cap\Vmin|$ and $\rho_{\max}:=|\mathrm{Rev}(G)\cap\Vmax|$. If
\[
  j\;\ge\;\min\bigl(k,\ \rho_{\min}+1\bigr)
  \qquad\text{then}\qquad \min_{x\in R_j}x=w^{\ast}\ \text{ pointwise},
\]
and dually $\max_{x\in R_j}x=w^{\ast}$ once
$j\ge\min(m,\rho_{\max}+1)$. In particular $R_j(G)=\{w^{\ast}\}$ as soon as
$j\ge\max\bigl(\min(k,\rho_{\min}+1),\min(m,\rho_{\max}+1)\bigr)$, so the first
exact level is at most $\min(m,k)$ and at most $\max(\rho_{\min},\rho_{\max})+1$.
\end{theorem}

\begin{proof}
Take a level-$j$ point and let $S\subseteq\Vmin$ be either $\Vmin$ itself (in
the first case) or $\mathrm{Rev}(G)\cap\Vmin$ (in the second); in both cases
$|S|\le j$, and in the second case $|S|+1\le j$. Iterating the splitting rule
(ii) one coordinate of $S$ at a time gives, for every $\tau\in\{0,1\}^{S}$ with
$S_\tau:=\tau^{-1}(0)$ and $T_\tau:=\tau^{-1}(1)$,
\[
  \sum_{\tau}z^{S_\tau,T_\tau}=1,\qquad
  \sum_{\tau}X^{S_\tau,T_\tau}=x .
\]
Fix $\tau$ with $\lambda_\tau:=z^{S_\tau,T_\tau}>0$ and put
$x^{\tau}:=X^{S_\tau,T_\tau}/\lambda_\tau$; by (iii) $x^{\tau}\in Q(G)$ and by
(iv) $x^{\tau}(v)=x^{\tau}(v^{(\tau(v))})$ for every $v\in S$. (If
$\lambda_\tau=0$ then $0\le X^{S_\tau,T_\tau}\le0$, so that index contributes
nothing.)

In the first case $S=\Vmin$ already. In the second case let $w\in\Vmin\setminus
S$ be rigid; since $|S|+1\le j$, splitting once more at $w$ writes $x^{\tau}$
as a convex combination of a point of $F_w^{0}$ and a point of $F_w^{1}$, so
$x^{\tau}(w)=\min_ix^{\tau}(w^{(i)})$ by \Cref{cl:rigid-pin}. Extend $\tau$ to
a profile $\tau'$ on $\Vmin$ by choosing at each such $w$ an action attaining
that minimum.

Either way $x^{\tau}(v)=x^{\tau}(v^{(\tau'(v))})$ at every Min vertex, while
$x^{\tau}(v)\ge x^{\tau}(v^{(i)})$ at every Max vertex and the average rows
hold; that is, $x^{\tau}\ge T_{\tau'}x^{\tau}$, where $T_{\tau'}$ is the
Shapley operator of the one-player game $G_{\tau'}$ obtained by freezing Min to
$\tau'$, and $x^{\tau}$ agrees with the sink payoffs. $G_{\tau'}$ is stopping
(a trap of $G_{\tau'}$ is a trap of $G$, by \Cref{lem:trapchar}), so
$T_{\tau'}$ is monotone with unique fixed point $\val^{\tau'}$ and
$T^{n}_{\tau'}x^{\tau}$ decreases to it; hence $x^{\tau}\ge\val^{\tau'}$.
Finally $\val^{\tau'}=\max_\sigma\val_{\sigma\tau'}\ge
\max_\sigma\min_{\tau''}\val_{\sigma\tau''}=w^{\ast}$ pointwise. Therefore
$x=\sum_\tau\lambda_\tau x^{\tau}\ge w^{\ast}$, and $w^{\ast}\in R_j$ gives
equality. The dual argument, with $\Vmax$ in place of $\Vmin$ and the reversed
inequalities, proves the second half; it is also \Cref{cl:duality} applied to
the first half.
\end{proof}

\begin{corollary}[Duality covariance]\label{cl:duality}
Let $\bar G$ be the dual game of \Cref{lem:duality}. Then
$R_j(\bar G)=\{\mathbf 1-x:x\in R_j(G)\}$ off the sinks, so
$\max_{R_j(\bar G)}x(v)=1-\min_{R_j(G)}x(v)$ and conversely.
\end{corollary}

\begin{proof}
The map $x\mapsto\mathbf 1-x$ exchanges the Max and Min order rows, fixes the
average rows and the box, and exchanges the two sink values; on the homogenised
level it is $X^{S,T}\mapsto z^{S,T}\mathbf 1-X^{S,T}$, which preserves (i),
(ii) and (iv) and maps (iii) for $G$ onto (iii) for $\bar G$.
\end{proof}

\begin{proposition}[Disjoint unions multiply]\label{cl:union}
Let $G$ be the SSG obtained from disjoint $G_1,G_2$ and a fresh average vertex
$s\to(r_1,r_2)$ with $r_i$ a vertex of $G_i$. Then
$R_j(G)|_{V_i}=R_j(G_i)$ for $i=1,2$ and every $j$, and
$\max_{R_j(G)}x(s)=\frac12\bigl(\max_{R_j(G_1)}x(r_1)+\max_{R_j(G_2)}x(r_2)\bigr)$,
and likewise for the minimum.
\end{proposition}

\begin{proof}
Restricting a level-$j$ point of $G$ to the indices $(S,T)$ with
$S\cup T\subseteq C_i$ and to the coordinates $V_i$ gives a level-$j$ point of
$G_i$, which proves $\subseteq$. Conversely, given level-$j$ points of $G_1$
and $G_2$, define for $(S,T)$ with $S_i:=S\cap C_i$, $T_i:=T\cap C_i$
\[
  z^{S,T}:=z_1^{S_1,T_1}z_2^{S_2,T_2},\qquad
  X^{S,T}|_{V_1}:=z_2^{S_2,T_2}X_1^{S_1,T_1},\quad
  X^{S,T}|_{V_2}:=z_1^{S_1,T_1}X_2^{S_2,T_2},
\]
and $X^{S,T}(s):=\frac12(X^{S,T}(r_1)+X^{S,T}(r_2))$. Every row of (iii) is
homogeneous and involves only one side, so it holds after the scaling;
splitting and tightness are inherited coordinatewise. The last display is the
average row at $s$ together with the product structure.
\end{proof}

\begin{theorem}[Level one is not exact]\label{cl:gap-one}
Let $W_{14}$ be the one-player stopping SSG of \Cref{thm:lasserre-vacuous}
($\Vmax=\{v_1,v_2,v_3\}$, $N=14$, $w^{\ast}=(\frac12,\frac12,\frac{15}{32})$ on
$(v_1,v_2,v_3)$). Then
\[
  \max_{x\in R_1(W_{14})}x(v_1)\;=\;\tfrac35\;>\;\tfrac12\;=\;w^{\ast}(v_1),
  \qquad
  \max_{x\in R_2(W_{14})}x(v_1)=\min_{x\in R_2(W_{14})}x(v_1)=\tfrac12 .
\]
Dually, on the dual game $\overline{W_{14}}$ (a Min-only stopping SSG on $14$
vertices), $\min_{R_1}x(v_1)=\frac25<\frac12=w^{\ast}(v_1)$.
\end{theorem}

\begin{proof}
In the coordinates $u=(x(v_1),x(v_2),x(v_3))$ the harmonic normal form of
$W_{14}$ (printed in \Cref{thm:lasserre-vacuous}) makes $Q(W_{14})$ the set
\[
  u_1\ge u_2,\quad u_1\ge u_3,\quad u_2\ge\tfrac12,\quad u_2\ge\tfrac78u_1,
  \quad u_3\ge\tfrac{15}{32},\quad u_3\ge\tfrac78u_1,\quad 0\le u\le1 ,
\]
with $F_{v_1}^{0}=\{u_1=u_2\}$, $F_{v_1}^{1}=\{u_1=u_3\}$,
$F_{v_2}^{0}=\{u_2=\frac12\}$, $F_{v_2}^{1}=\{u_2=\frac78u_1\}$,
$F_{v_3}^{0}=\{u_3=\frac{15}{32}\}$, $F_{v_3}^{1}=\{u_3=\frac78u_1\}$, all
intersected with $Q$. Put
\[
  u:=\bigl(\tfrac35,\tfrac{23}{40},\tfrac{11}{20}\bigr),\quad
  \pi:=\bigl(\tfrac8{13},\tfrac8{13},\tfrac7{13}\bigr),\quad
  \kappa:=\bigl(\tfrac47,\tfrac12,\tfrac47\bigr),\quad
  \pi':=\bigl(\tfrac12,\tfrac12,\tfrac{15}{32}\bigr),\quad
  \kappa':=\bigl(1,\tfrac78,\tfrac78\bigr).
\]
All five lie in $Q$: for $\pi$, $u_1=u_2$ and $u_3=\frac78u_1$;
for $\kappa$, $u_1=u_3$ and $u_2=\frac12=\frac78u_1$; $\pi'=w^{\ast}$;
$\kappa'$ has $u_2=u_3=\frac78u_1$; the remaining inequalities are the
displayed numerical checks $\frac8{13}\ge\frac12$, $\frac7{13}\ge\frac{15}{32}$,
$\frac47\ge\frac12$, $1\ge\frac78$. Moreover
\[
  u=\tfrac{13}{20}\pi+\tfrac7{20}\kappa
   =\tfrac45\pi'+\tfrac15\kappa' ,
\]
both identities being exact. The first exhibits $u\in\operatorname{conv}
(F_{v_1}^{0}\cup F_{v_1}^{1})$ (as $\pi\in F_{v_1}^{0}$,
$\kappa\in F_{v_1}^{1}$); the second exhibits
$u\in\operatorname{conv}(F_{v_2}^{0}\cup F_{v_2}^{1})$ and
$u\in\operatorname{conv}(F_{v_3}^{0}\cup F_{v_3}^{1})$ simultaneously, since
$\pi'\in F_{v_2}^{0}\cap F_{v_3}^{0}$ and
$\kappa'\in F_{v_2}^{1}\cap F_{v_3}^{1}$. By \Cref{cl:faces},
$u\in R_1(W_{14})$, so $\max_{R_1}x(v_1)\ge\frac35$. That the maximum is
exactly $\frac35$, and that level two is exact, were computed in exact rational
arithmetic. The dual statement is \Cref{cl:duality}.
\end{proof}

\begin{theorem}[The level-one relaxation does not decide the threshold]
\label{cl:dw}
Let $DW$ be the stopping SSG on $27$ vertices ($3$ Max, $3$ Min, $21$ average
including the sinks) obtained by joining $W_{14}$ and $\overline{W_{14}}$ under
a fresh average root $s\to(v_1,\bar v_1)$. Then
\[
  w^{\ast}(s)=\tfrac12,\qquad
  \min_{R_1}x(s)=\tfrac9{20}<\tfrac12<\tfrac{11}{20}=\max_{R_1}x(s),
  \qquad R_2(DW)=\{w^{\ast}\}.
\]
In particular the level-one relaxation determines neither $w^{\ast}(s)$ nor the
answer to \Cref{prob:main} at $s$, and no fixed level below two does.
\end{theorem}

\begin{proof}
$DW$ is stopping because $W_{14}$ and its dual are. The two displayed bounds
are \Cref{cl:union} applied to \Cref{cl:gap-one} and its dual:
$\frac12(\frac35+\frac12)=\frac{11}{20}$ and
$\frac12(\frac12+\frac25)=\frac9{20}$, while
$w^{\ast}(s)=\frac12(\frac12+\frac12)=\frac12$. Exactness at level two follows
from \Cref{cl:union} and level-two exactness on the two halves, and was also
computed directly on $DW$ in exact rational arithmetic.
\end{proof}

\begin{proposition}[The hierarchy on the instances of this paper]
\label{cl:families}
In exact rational arithmetic, $R_1=\{w^{\ast}\}$ --- the level-one relaxation
is the single point $w^{\ast}$, so one linear programme returns the whole value
vector --- on every one of
$G8$ (\Cref{prop:simorder-stalls}), $S$ (\Cref{prop:transport-stalls}),
$S_r$ with $r=3$ (\Cref{thm:transport-barrier}), $H_m$ for $m=3,5$
(\Cref{cor:slack-stalls}), the ten-vertex own-successor stall $R$
(\Cref{prop:own-stall}), $WD(e,j,m)$ for
$(e,j,m)=(4,2,6),(6,3,7),(8,4,8)$ (\Cref{def:wedge}), $CC(L,m)$ for
$(L,m)=(1,2),(2,2),(1,4)$ (\Cref{thm:hybrid-lower}), $CV(e,s)$ for
$(e,s)=(1,3),(1,4),(2,4)$ (\Cref{def:cv}), and $G^{\ast}$
(\Cref{prop:simorder-stalls}, $|C|=7$, $(m,k)=(3,4)$). On each of them
$R_0=Q(G)$ is \emph{not} a point.

Three of these are theorems rather than computations. $G8$, $H_m$ and $R$ have
$\mathrm{Rev}=\emptyset$ --- every controlled vertex is rigid --- so \Cref{cl:rigid}
applies. $WD$, $CC$, $S$, $S_r$ and every one-player family have
$\min(m,k)=0$, and $R$ has $\min(m,k)=1$, so \Cref{cl:collapse} applies.
$G^{\ast}$ has $\mathrm{Rev}(G^{\ast})=\{3,5,6\}\subseteq\Vmin$, hence
$\rho_{\max}=0$, so \Cref{cl:collapse} proves its upper half exact at level
one; only its lower half is a computation. Of the larger named instances,
$G^{\sharp}$ (\Cref{prop:auso-seven}, $(m,k)=(4,2)$) satisfies
$R_2=\{w^{\ast}\}$ and the level-two realisation of \Cref{prop:b2-realised}
($(m,k)=(5,1)$) satisfies $R_1=\{w^{\ast}\}$, by \Cref{cl:collapse} alone.
\end{proposition}

\begin{corollary}[What an instance defeating the hierarchy must contain]
\label{cl:barrier}
If $\max_{R_j}x(v)>w^{\ast}(v)$ for some vertex $v$ of a stopping $G$, then
$G$ has more than $j$ Max vertices and more than $j$ reversing Max vertices;
if in addition $\min_{R_j}x(v)<w^{\ast}(v)$ --- which is what it takes for the
level-$j$ bounds to fail to decide $\val(v)\ge\frac12$ --- then $G$ has more
than $j$ Min vertices and more than $j$ reversing Min vertices as well. In
particular no one-player game and no game with $|C|\le2$ can defeat
$R_1$ in both directions, and the entire negative arsenal of
\Cref{sec:gap} and \Cref{sec:wedge} --- every family there is one-player or has
$|C|\le2$ --- is decided by one linear programme of size $O(N^{2})$.
\end{corollary}

\begin{proof}
\Cref{cl:collapse} and \Cref{cl:sound}.
\end{proof}

\begin{corollary}[A cheap by-product: the gate cut]\label{cl:cut}
For $v\in C$ and $i\in\{0,1\}$ let
$M_{v,i}:=\max_{x\in Q(G)}\bigl|x(v)-x(v^{(i)})\bigr|$, computed by one linear
programme each. Then every $x\in R_1(G)$, and in particular $w^{\ast}$,
satisfies
\[
  \frac{|x(v)-x(v^{(0)})|}{M_{v,0}}+\frac{|x(v)-x(v^{(1)})|}{M_{v,1}}\ \le\ 1
  \qquad(v\in C),
\]
with the convention that a term with $M_{v,i}=0$ is read as $0$. These $|C|$
inequalities are valid for $w^{\ast}$, are not difference bounds in the sense
of \Cref{def:hybrid}, and cut $Q(G)$ strictly whenever some vertex is
value-distinguishing with both faces nonempty.
\end{corollary}

\begin{proof}
Write $h_i(x):=x(v^{(i)})-x(v)$; each $h_i$ has constant sign on $Q(G)$ and
$M_{v,i}=\max_{Q(G)}|h_i|$. With $x=\lambda\pi+(1-\lambda)\kappa$ as in
\Cref{cl:faces}, $h_0(\pi)=0$ gives $|h_0(x)|=(1-\lambda)|h_0(\kappa)|\le
(1-\lambda)M_{v,0}$ and $h_1(\kappa)=0$ gives
$|h_1(x)|\le\lambda M_{v,1}$; add.
\end{proof}

\begin{remark}[Where this sits among the paper's convex objects]
\label{cl:relation}
Four comparisons.
(i) $R_0=Q(G)$ exactly, so the hierarchy \emph{starts} at \Cref{def:transport};
\Cref{thm:transport-objective} says the whole difficulty there is the choice of
one of $2^{|C|}$ objectives, and \Cref{cl:rlt} replaces that choice by
$N^{O(j)}$ extra variables. \Cref{cl:rigid} identifies a polynomially
recognisable hypothesis --- no controlled vertex reverses its preference
anywhere on $Q(G)$ --- under which one further level already removes the
ambiguity.
(ii) The degree-two Lasserre lift of \Cref{def:lasserre-two} is a lift in the
\emph{value} variables, with the complementarity products $q_v$ and a
positive-semidefinite moment matrix; \Cref{cl:rlt} is a linear lift in the
\emph{choice} variables with no semidefinite constraint. On the two witnesses
of \Cref{thm:lasserre-vacuous} the comparison is decisive and in the opposite
direction to the one that theorem might suggest: on $W_{10}$ and $W_{14}$ the
degree-two lift returns the transport bound $1$ for $\max x(v_1)$, while
$\min_{R_0}x(v_1)=\frac12=w^{\ast}(v_1)$ already --- the \emph{minimum} over
the transport polytope itself is exact on both, which for $W_{14}$ is
\Cref{cl:collapse} with $k=0$ (equivalently the linear programme of
\Cref{thm:one-player}) and for $W_{10}$ is a computation, \Cref{cl:collapse}
giving level one there.
So \Cref{thm:lasserre-vacuous} is a statement about a direction, not about the
weakness of lifting; and $R_1$ improves $\max x(v_1)$ from $1$ to $\frac35$ on
$W_{14}$, $R_2$ to $\frac12$.
(iii) \Cref{sec:hybrid} cuts $Q(G)$ by \emph{difference} bounds
$z(p)-z(q)\le\Delta(p,q)$; $R_j$ cuts it by the projection of a lift, and
\Cref{cl:cut} exhibits a valid inequality of the latter that is not of the
former shape. Concretely, on $CC(L,m)$ the hybrid keeps
$\Delta_k(B_i,A_i)\ge0$ for $2^{\Omega(N)}$ rounds
(\Cref{thm:hybrid-lower}), the $Z$-seeded own-successor hybrid is silent for
$2^{\Omega(N)}$ rounds on $WD$ (\Cref{thm:wedge-proved},
\Cref{cor:wedge-count}) and on $CV$ (\Cref{thm:cv-lower}); one linear
programme,
$R_1$, returns the whole value vector on all three families
(\Cref{cl:families}).
(iv) Against \Cref{thm:modulator}: $\min(m,k)$ is not a modulator number, and
the bound $N^{O(\min(m,k))}$ of \Cref{cl:collapse} is dominated by the
$2^{\min(m,k)}\mathrm{poly}(N)$ enumeration noted at \Cref{thm:bsi-tracks};
what \Cref{cl:collapse} contributes is not a new complexity bound but the
\emph{level} at which this hierarchy collapses, and with it
\Cref{cl:barrier}, which says what an instance defeating it must contain.
The refined bound $\rho+1$ is genuinely smaller than both: $G^{\ast}$ has
$m=3$, $k=4$ and $\rho_{\max}=0$.
\end{remark}

\begin{remark}[The exact gap]\label{cl:gap}
What is settled: $R_0=Q(G)$, $R_1$ is not exact ($W_{14}$, $DW$), $R_j$ is
exact for $j\ge\min(m,k)$ and for $j\ge\max(\rho_{\min},\rho_{\max})+1$. What is
not settled is the single statement
\begin{quote}
\emph{there is a stopping simple stochastic game $G$ and a vertex $v$ with
$\max_{x\in R_2(G)}x(v)>w^{\ast}(v)$.}
\end{quote}
By \Cref{cl:barrier} such a $G$ has at least three Max vertices and at least
three reversing ones, and for the level-two \emph{bounds} to fail to decide
$\val(v)\ge\frac12$ it must have at least three reversing Min vertices as well.
Every instance named in this paper is excluded: those with
$\min(m,k)\le2$ or $\rho_{\max}\le2$ by \Cref{cl:collapse} --- which covers all
the families of \Cref{sec:gap} and \Cref{sec:wedge}, $G^{\ast}$ ($\rho_{\max}=0$),
$G^{\sharp}$ ($\min(m,k)=2$) and the level-two realisation of
\Cref{prop:b2-realised} ($\min(m,k)=1$) --- and the remaining ones,
$W_{14}$, $CV(e,s)$ and $DW$, by the computations of \Cref{cl:families},
\Cref{cl:gap-one} and \Cref{cl:dw}; and by \Cref{cl:union} no disjoint union of
instances that fail only at level one can serve either. The negation of the
displayed statement --- that $\max_{R_2}x$ is $w^{\ast}$ on every stopping game
--- would put \Cref{prob:main} in $\mathsf P$, the optimum of a rational linear
programme of size $N^{O(1)}$ being computable exactly; it is therefore
target-equivalent and is not conjectured here.
\end{remark}
